\documentclass[10pt, a4paper, landscape]{article}

% -------------------------------------------------
% 宏包引入
% -------------------------------------------------
\usepackage[fontset=mac]{ctex}       % 中文支持
\usepackage{multicol}   % 多分栏
\usepackage{calc}
\usepackage{ifthen}
\usepackage[landscape]{geometry} % 页面设置
\usepackage{amsmath,amsthm,amsfonts,amssymb} % 数学公式
\usepackage{mathtools} % 提供 \xleftrightarrow 等箭头
\usepackage{color,graphicx,overpic} % 颜色与图片
\usepackage{hyperref}   % 超链接
\usepackage{enumitem}   % 列表环境控制
\usepackage{titlesec}   % 标题控制
\usepackage{bm}         % 加粗数学符号
\usepackage{xcolor}
\usepackage{tabularx}
\usepackage{array}
\usepackage{booktabs}
\usepackage{color} % 或 xcolor
\usepackage{tikz}       % 绘图
\usetikzlibrary{decorations.pathreplacing, positioning} % 加载brace装饰库
\usetikzlibrary{calc, positioning, arrows.meta, decorations.markings}
\setCJKmainfont{PingFang SC}
\setCJKsansfont{PingFang SC}
\setCJKmonofont{PingFang SC}

% -------------------------------------------------
% 自定义颜色
% -------------------------------------------------

\definecolor{myblue}{HTML}{003153} % 蓝色字
\definecolor{myred}{HTML}{85120F}  % 红色字
\definecolor{mygreen}{HTML}{218254}  % 绿色字
\definecolor{hlblue}{HTML}{ADC1DB} % 蓝色高亮
\definecolor{hlred}{HTML}{740003}  % 红色高亮
\definecolor{hlred40}{HTML}{D5807F}  % 红色高亮40%透明度
\definecolor{sysugreen}{HTML}{00561F}  
\definecolor{hlyellow}{HTML}{E3C79F}  % 奢金高亮
\definecolor{hlgreen}{HTML}{BED49D}  % 抹茶绿高亮


% -------------------------------------------------
% 极限空间压缩设置 (核心部分)
% -------------------------------------------------

% 1. 页边距设置为极小 (0.5cm)
\geometry{top=0.5cm,left=0.5cm,right=0.5cm,bottom=0.5cm}

% 2. 去掉段落首行缩进，改为段落间略微留空（可选，这里为了紧凑设为0）
\setlength{\parindent}{0pt}
\setlength{\parskip}{0pt}

% 3. 设置正文基础字体大小为 scriptsize (约8pt)，如果还觉得大，可以改为 \tiny
\renewcommand{\baselinestretch}{0.9} % 压缩行间距
\let\oldfootnotesize\footnotesize
\renewcommand{\footnotesize}{\fontsize{7pt}{8pt}\selectfont}

% 4. 压缩列表环境 (Itemize/Enumerate) 的间距
\setlist{nolistsep} 
\setlist[itemize]{leftmargin=*}
\setlist[enumerate]{leftmargin=*}

% 5. 压缩标题间距
\titleformat{\section}{\bfseries\scriptsize\color{myblue}}{}{0em}{}[\hrule] % 标题带下划线，蓝色，省空间
\titlespacing*{\section}{0pt}{2pt}{1pt} % 上方留2pt，下方留1pt
\titleformat{\subsection}
    [runin] % 不换行
    {\bfseries\scriptsize} % 粗体、scriptsize，黑色字体
    {} % 不显示编号
    {0pt} % 标题与正文间距
    {\hlyellow} % 用hlyellow高亮命令包裹标题
    [] % 标题内容后无内容
\titleformat{\subsubsection}
    [runin] % 不换行
    {\bfseries\tiny} % 粗体、scriptsize，黑色字体
    {} % 不显示编号
    {0pt} % 标题与正文间距
    {\hlgreen} % 用hlgreen高亮命令包裹标题
    [] % 标题内容后无内容
\titlespacing*{\subsection}{0pt}{1pt}{0.5em} % 上方1pt,下方0.5em(水平间距)
\titlespacing*{\subsubsection}{0pt}{1pt}{0.5em} % 上方1pt,下方0.5em(水平间距)

% -------------------------------------------------
% 自定义命令
% -------------------------------------------------
% 颜色字
\newcommand{\red}[1]{\textbf{\textcolor{myred}{#1}}}  
\newcommand{\blue}[1]{\textbf{\textcolor{myblue}{#1}}} 
\newcommand{\green}[1]{\textbf{\textcolor{mygreen}{#1}}}
\newcommand{\entry}[2]{$\bullet$ \textbf{#1}: #2\par\vspace{0.5pt}}
% 高亮
\newcommand{\cbox}[2][yellow]{\begingroup\setlength{\fboxsep}{1pt}\colorbox{#1}{\strut#2}\endgroup}
\newcommand{\hlblue}[1]{\cbox[hlblue]{#1}}
\newcommand{\hlred}[1]{\cbox[hlred40]{#1}}
\newcommand{\hlyellow}[1]{\cbox[hlyellow]{#1}}
\newcommand{\hlgreen}[1]{\cbox[hlgreen]{#1}} 
% 图片插入
\newcommand{\img}[2][0.9\linewidth]{%
    {\par\vspace{1pt}\centering\includegraphics[width=#1]{#2}\par\vspace{1pt}}%
}
% 左图右文 (参数: [图片宽度比例]{图片路径}{右侧文字内容})
\newcommand{\imgleft}[3][0.3]{%
    \noindent\begin{minipage}[t]{#1\linewidth}%
        \vspace{0pt}%
        \includegraphics[width=\linewidth]{#2}%
    \end{minipage}%
    \hfill%
    \begin{minipage}[t]{0.98\linewidth - #1\linewidth}%
        \vspace{0pt}%
        #3%
    \end{minipage}\par\vspace{2pt}%
}
% 左文右图 (参数: [图片宽度比例]{图片路径}{左侧文字内容})
\newcommand{\imgright}[3][0.3]{%
    \noindent\begin{minipage}[t]{0.98\linewidth - #1\linewidth}%
        \vspace{0pt}%
        #3%
    \end{minipage}%
    \hfill%
    \begin{minipage}[t]{#1\linewidth}%
        \vspace{0pt}%
        \includegraphics[width=\linewidth]{#2}%
    \end{minipage}\par\vspace{2pt}%
}
% -------------------------------------------------
% 新增：概念速查表专用命令
% -------------------------------------------------
% 表格容器
\newcommand{\concepttable}[1]{%
    {\setlength{\tabcolsep}{1.5pt}% 局部减小列间距
     \renewcommand{\arraystretch}{0.92}% 局部紧缩行距
     \par\vspace{2pt}{\color{myblue}\hrule height 0.6pt}\vspace{1pt}% 上边框(蓝色,0.6pt粗)
     \noindent\begin{tabular}{@{}p{0.22\linewidth}p{0.76\linewidth}@{}}%
     #1%
     \end{tabular}%
     \vspace{1pt}{\color{myblue}\hrule height 0.6pt}\par\vspace{2pt}}% 下边框(蓝色,0.6pt粗)
}
% 表格行 (参数: {概念名}{解释})
\newcommand{\conrow}[2]{\blue{#1} & #2 \\}


% 定义基于 X 列的自适应等宽列：L (左对齐) 和 C (居中对齐)
\newcolumntype{L}{>{\raggedright\arraybackslash}X}
\newcolumntype{C}{>{\centering\arraybackslash}X}

% ==========================================
% 智能满宽、均匀等分布对比表格环境 \compTable
% 参数说明：
%   [#1]: 可选参数，右侧对比列的格式，默认 C (居中且等宽)
%   {#2}: 必填参数，表格的总列数 (数字)
%   {#3}: 必填参数，表格的完整内容 (包括 header 和 row)
% ==========================================
\newcommand{\compTable}[3][C]{%
    {\setlength{\tabcolsep}{6pt}% 适度保持列间距
     \renewcommand{\arraystretch}{1.15}% 稍微舒展一点的行距
     \par\vspace{3pt}{\color{myblue}\hrule height 0.8pt}\vspace{3pt}% 上边框
     \noindent
     \def\colsnum{#2}%
     % 核心改动：改用 tabularx，宽度锁定为 \linewidth，使用 L 和 C 强制均匀平分列宽
     \begin{tabularx}{\linewidth}{@{} L *{\numexpr\colsnum-1\relax}{#1} @{}}%
         #3%
     \end{tabularx}%
     \vspace{2pt}{\color{myblue}\hrule height 0.8pt}\par\vspace{3pt}}% 下边框
}

% 泛用表格行 \comprow
\newcommand{\comprow}[1]{#1 \\}

% 专门的表头行 \compheader (自带下划线，不使用 \textbf 避免报错)
\newcommand{\compheader}[1]{%
    #1 \\ \noalign{\vspace{1.5pt}{\color{myblue}\hrule height 0.4pt}\vspace{2.5pt}}%
}

% 题干专属方框 (minipage版)
% 用法: \prob{大题号}{小题号}{题干内容}，小题号可空
% 框宽占满行宽，超出行宽时自动换行
\newcommand{\prob}[3]{%
    \par\vspace{3pt}\noindent%
    \fcolorbox{myblue}{myblue!5}{%
        \begin{minipage}{\dimexpr\linewidth-2\fboxsep-2\fboxrule\relax}%
        \vspace{1pt}%
        \textbf{\textcolor{myblue}{[T#1\if\relax\detokenize{#2}\relax\else(#2)\fi]}} #3%
        \vspace{1pt}%
        \end{minipage}%
    }%
    \par\vspace{3pt}%
}

% 注释专属方框 (minipage版)
% 用法: \comment{大题号}{小题号}{注释内容}，小题号可空
% 框宽占满行宽，超出行宽时自动换行
\newcommand{\comment}[3]{%
    \par\vspace{3pt}\noindent%
    \fcolorbox{sysugreen}{sysugreen!5}{%
        \begin{minipage}{\dimexpr\linewidth-2\fboxsep-2\fboxrule\relax}%
        \vspace{1pt}%
        \textbf{\textcolor{sysugreen}{[注#1\if\relax\detokenize{#2}\relax\else(#2)\fi]}} #3%
        \vspace{1pt}%
        \end{minipage}%
    }%
    \par\vspace{3pt}%
}

% -------------------------------------------------
% 正文
% -------------------------------------------------
\begin{document}

\tiny

% 三栏布局
\begin{multicols*}{3}

\section{第一章\ 采样定理}

\section{第二章\ 时域离散信号及系统频域分析}

\prob{1}{}{设$X(e^{j\omega})$和$Y(e^{j\omega})$分别是$x(n)$和$y(n)$的傅里叶变换，试求下面序列的傅里叶变换：（1）$x(n-n_0)$；（2）$x^*(n)$；（3）$x(-n)$；（4）$x(n)*y(n)$；（5）$x(n)y(n)$；（6）$nx(n)$；（7）$x(2n)$；（8）$x^2(n)$；（9）$x_9(n)=\begin{cases}x\left(\frac{n}{2}\right) & n=\text{偶数} \\ 0 & n=\text{奇数}\end{cases}$。}

解：

(1)～(6)都可以通过查性质表得到。

(7)注意$x(2n) \stackrel{\mathcal{F}}{\longleftrightarrow} \frac{1}{2}\left[X\left(\mathrm{e}^{\mathrm{j}\frac{\omega}{2}}\right) + X\left(\mathrm{e}^{\mathrm{j}\left(\frac{\omega}{2}-\pi\right)}\right)\right]$

(8)x(n)和自己相乘，即应用 DTFT 的时域相乘性质即可。

(9)略

\prob{\textbf{\hlblue{2}}}{}{已知$X(e^{j\omega}) = \begin{cases} 1, & |\omega| < \omega_0 \\ 0, & \omega_0 < |\omega| \leq \pi \end{cases}$，求$X(e^{j\omega})$的傅里叶反变换$x(n)$。}

\textbf{解：}

由 IDTFT 定义：
\[
x(n) = \frac{1}{2\pi} \int_{-\pi}^{\pi} X(e^{j\omega})\,e^{j\omega n}\,d\omega
\]

$\because\ X(e^{j\omega})$ 仅在 $|\omega|<\omega_0$ 时为 $1$，其余为零，

$\therefore$ 积分限可收缩至 $[-\omega_0,\ \omega_0]$：
\[
x(n) = \frac{1}{2\pi} \int_{-\omega_0}^{\omega_0} e^{j\omega n}\,d\omega
     = \frac{\sin(\omega_0 n)}{n\pi}
\]

\textbf{\textcolor{mygreen}{注：$X(e^{j\omega})$ 为经典的矩形窗频谱，其 IDTFT 为 sinc 函数形式，需熟记。}}

\prob{3}{}{线性时不变系统的频率响应(频率响应函数)$H(e^{j\omega}) = |H(e^{j\omega})|e^{j\theta(\omega)}$，如果单位脉冲响应$h(n)$为实序列，试证明输入$x(n) = A\cos(\omega_0 n + \varphi)$的稳态响应为$y(n) = A|H(e^{j\omega_0})|\cos\left[\omega_0 n + \varphi + \theta(\omega_0)\right]$}

\textbf{解：}

实际上就是课件中的一个性质的证明（\textbf{单频率周期信号输入，系统输出为单频率周期信号，且幅值和相位对应变化}），当成结论记住即可。

\prob{4}{}{设$x(n) = \begin{cases} 1, & n = 0,1 \\ 0, & \text{其他} \end{cases}$，将$x(n)$以4为周期进行周期延拓，形成周期序列$\widetilde{x}(n)$，画出$x(n)$和$\widetilde{x}(n)$的波形，求出$\widetilde{x}(n)$的离散傅里叶级数$\widetilde{X}(k)$和傅里叶变换。}

\textbf{解：}

\vspace{4pt}
\begin{itemize}[leftmargin=*,itemsep=6pt]
\item[\textbf{1)}] \textbf{求 DFS：$\widetilde{X}(k)$}

$\widetilde{x}(n)$ 由 $x(n)=\{1,1\}$ 以 $N=4$ 周期延拓得到。
\[
\widetilde{X}(k)=\mathrm{DFS}[\widetilde{x}(n)]
                 =\sum_{n=0}^{3}\widetilde{x}(n)e^{-j\frac{2\pi}{4}kn}
                 =\sum_{n=0}^{1}e^{-j\frac{\pi}{2}kn}
                 =1+e^{-j\frac{\pi}{2}k}
\]
提公因子化为幅度-相位形式：
\[
\widetilde{X}(k)=e^{-j\frac{\pi}{4}k}
                  \bigl(e^{j\frac{\pi}{4}k}+e^{-j\frac{\pi}{4}k}\bigr)
                =2\cos\!\Bigl(\frac{\pi}{4}k\Bigr)\,
                  e^{-j\frac{\pi}{4}k},
\qquad \widetilde{X}(k)\ \text{以 $4$ 为周期}.
\]

\item[\textbf{2)}] \textbf{求 FT：$X(e^{j\omega})$}

周期序列的 FT 为频域冲激串（$N=4$）：
\[
X(e^{j\omega})=\mathrm{FT}[\widetilde{x}(n)]
              =\frac{2\pi}{N}\sum_{k=-\infty}^{\infty}
                \widetilde{X}(k)\,
                \delta\!\Bigl(\omega-\frac{2\pi}{N}k\Bigr)
              =\frac{\pi}{2}\sum_{k=-\infty}^{\infty}
                \widetilde{X}(k)\,
                \delta\!\Bigl(\omega-\frac{\pi}{2}k\Bigr).
\]
代入 $\widetilde{X}(k)$ 表达式：
\[
\boxed{%
X(e^{j\omega})=\pi\sum_{k=-\infty}^{\infty}
               \cos\!\Bigl(\frac{\pi}{4}k\Bigr)
               e^{-j\frac{\pi}{4}k}\,
               \delta\!\Bigl(\omega-\frac{\pi}{2}k\Bigr)
}.
\]
\end{itemize}

\prob{\textbf{\hlblue{5}}}{}{设题5图所示的序列$x(n)$的FT用$X(e^{j\omega})$表示，不直接求出$X(e^{j\omega})$，完成下列运算或工作：\\
(1) $X(e^{j0})$；\\
(2) $\int_{-\pi}^{\pi} X(e^{j\omega}) d\omega$；\\
(3) $X(e^{j\pi})$；\\
(4) 确定并画出傅里叶变换实部$\text{Re}[X(e^{j\omega})]$的时间序列$x_o(n)$；\\
(5) $\int_{-\pi}^{\pi} \left| X(e^{j\omega}) \right|^2 d\omega$；\\
(6) $\int_{-\pi}^{\pi} \left| \frac{dX(e^{j\omega})}{d\omega} \right|^2 d\omega$。\\}

\textbf{解：}

\begin{itemize}[leftmargin=*,itemsep=4pt]
\item[(1)] \textbf{\textcolor{myred}{$X(e^{j0})$ — 直流分量}}
  \[
  X(e^{j0}) = \sum_{n=-\infty}^{\infty} x(n) e^{-j0\cdot n}
            = \sum_{n=-\infty}^{\infty} x(n)
            = \sum_{n=-3}^{7} x(n) = 6
  \]
  \textcolor{myred}{要点：$\omega=0 \ \Rightarrow\ e^{-j\omega n}=1$，即序列各点之和。}

\item[(2)] \textbf{\textcolor{myred}{$\displaystyle\int_{-\pi}^{\pi} X(e^{j\omega})\,\mathrm{d}\omega$ — 反变换求 $x(0)$}}
  \[
  x(n) = \frac{1}{2\pi}\int_{-\pi}^{\pi} X(e^{j\omega})\,e^{j\omega n}\,\mathrm{d}\omega
  \quad\overset{n=0}{\Longrightarrow}\quad
  x(0) = \frac{1}{2\pi}\int_{-\pi}^{\pi} X(e^{j\omega})\,\mathrm{d}\omega
  \]
  \[
  \therefore\ \int_{-\pi}^{\pi} X(e^{j\omega})\,\mathrm{d}\omega = 2\pi\cdot x(0) = 4\pi
  \]

\item[(3)] \textbf{\textcolor{myred}{$X(e^{j\pi})$ — Nyquist 频率}}
  \[
  X(e^{j\pi}) = \sum_{n=-\infty}^{\infty} x(n)\,e^{-j\pi n}
              = \sum_{n=-3}^{7} (-1)^n x(n) = 2
  \]
  \textcolor{myred}{要点：$e^{-j\pi n}=(-1)^n$，序列交错求和。}

\item[(4)] \textbf{\textcolor{myred}{$\operatorname{Re}[X(e^{j\omega})]$ 对应的时间序列 $x_e(n)$}}
  \[
  \boxed{x_e(n) = \tfrac{1}{2}\bigl[x(n) + x^*(-n)\bigr]}
  \xrightarrow{\text{实序列}}\
  x_e(n) = \tfrac{1}{2}\bigl[x(n) + x(-n)\bigr]
  \]
  \textcolor{myred}{即共轭对称分量，其傅里叶变换为 $\operatorname{Re}[X(e^{j\omega})]$.}

\item[(5)] \textbf{\textcolor{myred}{$\displaystyle\int_{-\pi}^{\pi} |X(e^{j\omega})|^2\,\mathrm{d}\omega$ — 帕塞瓦尔定理}}
  \[
  \int_{-\pi}^{\pi} |X(e^{j\omega})|^2\,\mathrm{d}\omega
  = 2\pi\!\!\sum_{n=-\infty}^{\infty} |x(n)|^2
  = 2\pi\!\!\sum_{n=-3}^{7} |x(n)|^2 = 28\pi
  \]

\item[(6)] \textbf{\textcolor{myred}{$\displaystyle\int_{-\pi}^{\pi} \bigl|\tfrac{\mathrm{d}X}{\mathrm{d}\omega}\bigr|^2\,\mathrm{d}\omega$ — 频域微分 + 帕塞瓦尔}}
  \[
  \frac{\mathrm{d}X(e^{j\omega})}{\mathrm{d}\omega}
  = \operatorname{FT}\bigl[-jn\,x(n)\bigr]
  \]
  \[
  \therefore\ \int_{-\pi}^{\pi} \Bigl|\frac{\mathrm{d}X}{\mathrm{d}\omega}\Bigr|^2\,\mathrm{d}\omega
  = 2\pi\!\!\sum_{n=-3}^{7} |n\,x(n)|^2 = 316\pi
  \]
\end{itemize}

\prob{6}{}{试求如下序列的傅里叶变换：\\
(1) $x_1(n) = \delta(n-3)$；\\
(2) $x_2(n) = \frac{1}{2}\delta(n+1) + \delta(n) + \frac{1}{2}\delta(n-1)$；\\
(3) $x_3(n) = a^n u(n)\ (0 < a < 1)$；\\
(4) $x_4(n) = u(n+3) - u(n-4)$。}

\textbf{解：}

(1) 由 $\delta(n-3)$ 的筛选性，仅有 $n=3$ 项非零：
  \[
  X_1(e^{j\omega}) = \sum_{n=-\infty}^{\infty} \delta(n-3) e^{-j\omega n} = e^{-j3\omega}
  \]

(2) 序列仅在 $n=-1,0,1$ 处非零，直接代入 FT 定义，并利用欧拉公式合并共轭项：
  \[
  X_2(e^{j\omega}) = \sum_{n=-\infty}^{\infty} x_2(n) e^{-j\omega n}
  = \frac{1}{2}e^{j\omega} + 1 + \frac{1}{2}e^{-j\omega}
  = 1 + \cos\omega
  \]

(3) 由 $u(n)$ 截断负半轴，使求和下限为 $n=0$；再利用几何级数公式 $\sum_{n=0}^{\infty} r^n = \frac{1}{1-r}$（$|r|<1$），此处 $r = a e^{-j\omega}$ 且 $0<a<1$ 保证收敛：
  \[
  X_3(e^{j\omega}) = \sum_{n=-\infty}^{\infty} a^n u(n) e^{-j\omega n}
  = \sum_{n=0}^{\infty} \bigl(a e^{-j\omega}\bigr)^n
  = \frac{1}{1 - a e^{-j\omega}}
  \]

(4) $u(n+3)-u(n-4)$ 为长度 7 的矩形窗（$n\in[-3,3]$），可分正负半轴求和，然后分别用有限项等比级数公式收取：
  \[
  \begin{aligned}
  X_4(e^{j\omega})
  &= \sum_{n=-\infty}^{\infty} \bigl[u(n+3) - u(n-4)\bigr] e^{-j\omega n}
   = \sum_{n=-3}^{3} e^{-j\omega n} \\[4pt]
  &= \sum_{n=0}^{3} e^{-j\omega n} + \sum_{n=-3}^{-1} e^{-j\omega n}
   = \frac{1 - e^{-j4\omega}}{1 - e^{-j\omega}} + \frac{e^{j\omega} - e^{-j3\omega}}{1 - e^{-j\omega}} \cdot e^{j\omega}
   = \frac{1 - e^{-j4\omega}}{1 - e^{-j\omega}} + \frac{e^{j\omega}\bigl(1 - e^{-j3\omega}\bigr)}{1 - e^{-j\omega}} e^{j\omega}
  \end{aligned}
  \]

\prob{7}{}{设：\\
(1) $x(n)$是实偶函数；\\
(2) $x(n)$是实奇函数；\\
分别分析推导以上两种假设下，其$x(n)$的傅里叶变换性质。}

\textbf{解：}

记住对应结论（证明的分不要了）：

实偶 $\rightarrow$ 实偶，实奇 $\rightarrow$ 虚奇\\
虚奇 $\rightarrow$ 实奇，虚偶 $\rightarrow$ 虚偶

\prob{8}{}{设$x(n) = R_4(n)$，试求$x(n)$的共轭对称序列$x_e(n)$和共轭反对称序列$x_o(n)$，并分别用图表示。}

\textbf{解：}

图略，直接按照奇偶函数的定义式求解即可。

\prob{9}{}{已知$x(n) = a^n u(n)\ (0 < a < 1)$，分别求出其偶函数$x_e(n)$和奇函数$x_o(n)$的傅里叶变换。}

\textbf{解：}

先求 $x(n)$ 的傅里叶变换：

\[
X(e^{j\omega}) = \frac{1}{1 - a e^{-j\omega}}.
\]

由共轭对称性质：

\textbf{\textcolor{myred}{\[
\text{FT}[x_e(n)] = \operatorname{Re}\bigl[X(e^{j\omega})\bigr], \qquad
\text{FT}[x_o(n)] = j \operatorname{Im}\bigl[X(e^{j\omega})\bigr].
\]}}

为提取实部和虚部，将 $X(e^{j\omega})$ 分母实数化：

\[
X(e^{j\omega}) = \frac{1}{1 - a e^{-j\omega}}
               \cdot \frac{1 - a e^{j\omega}}{1 - a e^{j\omega}}
               = \frac{1 - a e^{j\omega}}{1 + a^2 - 2a \cos\omega}.
\]

分离实部与虚部：

\[
\operatorname{Re}[X(e^{j\omega})]
 = \frac{1 - a \cos\omega}{1 + a^2 - 2a \cos\omega}, \qquad
\operatorname{Im}[X(e^{j\omega})]
 = \frac{-a \sin\omega}{1 + a^2 - 2a \cos\omega}.
\]

故：

\[
\boxed{\text{FT}[x_e(n)] = \frac{1 - a \cos\omega}{1 + a^2 - 2a \cos\omega}}, \qquad
\boxed{\text{FT}[x_o(n)] = \frac{-a \sin\omega}{1 + a^2 - 2a \cos\omega}}.
\]


\prob{10}{}{若序列$h(n)$是\textbf{\textcolor{myred}{实因果序列}}，其傅里叶变换的实部如下式：$H_R(e^{j\omega}) = 1 + \cos\omega$，求序列$h(n)$及其傅里叶变换$H(e^{j\omega})$。}

\comment{1}{}{
\textbf{\textcolor{myred}{常见函数的复指数形式：}}
\begin{align*}
\cos\omega = \frac{e^{j\omega} + e^{-j\omega}}{2}, &\quad \sin\omega = \frac{e^{j\omega} - e^{-j\omega}}{2j}
\end{align*}
}

\textbf{解：}

由傅里叶变换的性质，实部 $H_R(e^{j\omega})$ 对应于序列共轭对称分量 $h_e(n)$ 的傅里叶变换。\textbf{将已知式展为复指数形式}：

\[
H_R(e^{j\omega}) = 1 + \cos\omega = 1 + \tfrac{1}{2}e^{j\omega} + \tfrac{1}{2}e^{-j\omega} = \text{FT}[h_e(n)].
\]

比较 $\sum_{n} h_e(n) e^{-j\omega n}$，直接读出共轭对称分量：

\[
h_e(n) = \begin{cases}
\frac{1}{2}, & n = -1, \\
1,        & n = 0, \\
\frac{1}{2}, & n = 1, \\
0,        & \text{其它}.
\end{cases}
\]

\comment{2}{实因果序列与共轭（反）对称分量的关系}{
对于实因果序列 $x(n)$满足$x(n) = x_e(n) + x_o(n)$，其中：
\begin{itemize}[leftmargin=*]
  \item 共轭对称分量：$x_e(n) = \frac{1}{2}[x(n) + x^*(-n)]$
  \item 共轭反对称分量：$x_o(n) = \frac{1}{2}[x(n) - x^*(-n)]$
\end{itemize}

因果性使二者不是独立的，可由其一恢复 $x(n)$：
\[
x(n) = \begin{cases}
0,              & n < 0, \\[2pt]
x_e(0),         & n = 0, \\[2pt]
2x_e(n),        & n > 0.
\end{cases}
\qquad
x(n) = \begin{cases}
0,              & n < 0, \\[2pt]
x(0),           & n = 0, \\[2pt]
2x_o(n),        & n > 0.
\end{cases}
\]
}

代入得：

\[
\boxed{h(n) = \begin{cases}
1, & n = 0, \\
1, & n = 1, \\
0, & \text{其它}.
\end{cases}}
\]

对其直接作离散时间傅里叶变换：

\[
H(e^{j\omega}) = \sum_{n=-\infty}^{\infty} h(n) e^{-j\omega n}
               = 1 + e^{-j\omega}
               = 2e^{-j\omega/2} \cos\frac{\omega}{2}.
\]

\[
\boxed{H(e^{j\omega}) = 2e^{-j\omega/2} \cos(\omega/2)}.
\]

\prob{11}{}{若序列$h(n)$是实因果序列，$h(0) = 1$，其傅里叶变换的虚部为$H_I(e^{j\omega}) = -\sin\omega$，求序列$h(n)$及其傅里叶变换$H(e^{j\omega})$。}

\textbf{解：}

由已知，傅里叶变换的虚部为
\[
H_I(e^{j\omega}) = -\sin\omega = -\frac{1}{2j}\bigl[e^{j\omega} - e^{-j\omega}\bigr].
\]

对于实序列，\textbf{\textcolor{myred}{其奇部 $h_o(n)$ 的傅里叶变换满足 $\text{FT}[h_o(n)] = jH_I(e^{j\omega})$}}，故
\[
\text{FT}[h_o(n)] = jH_I(e^{j\omega}) 
= -\frac{1}{2}\bigl[e^{j\omega} - e^{-j\omega}\bigr]
= \sum_{n=-\infty}^{\infty} h_o(n)\,e^{-j\omega n}.
\]

对照逆变换系数，立得奇部序列：
\[
h_o(n) = \begin{cases}
-\dfrac{1}{2}, & n = -1, \\[4pt]
0,            & n = 0,  \\[4pt]
\dfrac{1}{2}, & n = 1.
\end{cases}
\]

对于实因果序列，有还原关系
\[
h(n) = \begin{cases}
0,            & n < 0, \\
h(0) = 1,     & n = 0, \\
2h_o(n),      & n > 0.
\end{cases}
\]

由 $h(0)=1$ 及 $h_o(1)=\frac{1}{2}$ 得 $h(1)=2h_o(1)=1$，其余 $n>0$ 处 $h_o(n)=0$，故
\[
\boxed{h(n) = \begin{cases}
1, & n = 0, \\
1, & n = 1, \\
0, & \text{其它}.
\end{cases}}
\]

对其直接作离散时间傅里叶变换：
\[
H(e^{j\omega}) = \sum_{n=-\infty}^{\infty} h(n)\,e^{-j\omega n}
               = 1 + e^{-j\omega}
               = 2e^{-j\omega/2}\cos\frac{\omega}{2}.
\]

\[
\boxed{H(e^{j\omega}) = 2e^{-j\omega/2}\cos(\omega/2)}.
\]

\prob{12}{}{设系统的单位脉冲响应$h(n)=a^n u(n)$，$0<a<1$，输入序列为$x(n)=\delta(n)+2\delta(n-2)$，完成下面各题：\\ (1) 求出系统输出序列$y(n)$；\\ (2) 分别求出$x(n)$、$h(n)$和$y(n)$的傅里叶变换。}

\textbf{解：}

\textbf{(1) 求输出序列 $y(n)$。}
由卷积运算及单位脉冲的筛选性质，有：
\[
y(n) = h(n) * x(n)
     = a^n u(n) * \bigl[\delta(n) + 2\delta(n-2)\bigr]
     = a^n u(n) + 2a^{n-2} u(n-2).
\]
其中第二项利用了 $\delta(n-2)$ 将 $h(n)$ 平移两位并乘以 2 的性质。\\[4pt]
\textbf{(2) 求 $x(n)$、$h(n)$、$y(n)$ 的傅里叶变换。}
不宜直接套用定义式逐项求和，而应借助经典傅里叶变换对及性质巧解：

$\bullet$ $x(n)$：利用 $\delta(n) \xleftrightarrow{\mathcal{F}} 1$ 及时移性质 $\delta(n-n_0) \xleftrightarrow{\mathcal{F}} e^{-j\omega n_0}$，得
\[
X(e^{j\omega}) = 1 + 2e^{-j2\omega}.
\]

$\bullet$ $h(n)$：由经典变换对 $a^n u(n) \xleftrightarrow{\mathcal{F}} \dfrac{1}{1-ae^{-j\omega}}$（$|a|<1$），直接写出
\[
H(e^{j\omega}) = \frac{1}{1 - a e^{-j\omega}}.
\]

$\bullet$ $y(n)$：由傅里叶变换的卷积定理 $y(n)=x(n)*h(n) \;\xleftrightarrow{\mathcal{F}}\; Y(e^{j\omega})=X(e^{j\omega})H(e^{j\omega})$，得
\[
Y(e^{j\omega}) = H(e^{j\omega})\,X(e^{j\omega})
              = \frac{1 + 2e^{-j2\omega}}{1 - a e^{-j\omega}}.
\]

\prob{14}{}{求出以下序列的Z变换及收敛域：\\ (1) $2^{-n} u(n)$ \quad (2) $-2^{-n} u(-n-1)$ \\ (3) $2^{-n} u(-n)$ \quad (4) $\delta(n)$ \\ (5) $\delta(n-1)$ \quad (6) $2^{-n}[u(n)-u(n-10)]$}

\textbf{解：} 

\textbf{(1) $2^{-n}u(n)$}：令 $a=2^{-1}$，直接由经典变换对得
\[
\mathcal{Z}\big[2^{-n}u(n)\big] = \frac{1}{1-\frac{1}{2}z^{-1}},\qquad |z|>\frac{1}{2}.
\]

\textbf{(2) $-2^{-n}u(-n-1)$}：此为左边序列。由经典左边序列变换对
\[
-a^{n}u(-n-1) \xleftrightarrow{Z} \frac{1}{1-az^{-1}},\;|z|<|a|,
\]
取 $a=2^{-1}$ 即得
\[
\mathcal{Z}\big[-2^{-n}u(-n-1)\big] = \frac{1}{1-\frac{1}{2}z^{-1}},\qquad |z|<\frac{1}{2}.
\]

\textbf{\textcolor{mygreen}{注意分子分母表达式与(1)相同，但收敛域截然不同}}

\textbf{(3) $2^{-n}u(-n)$}：$u(-n)$ 在 $n\le 0$ 时为1，令 $m=-n$，化为右边序列的幂级数：
\[
\mathcal{Z}\big[2^{-n}u(-n)\big] = \sum_{n=-\infty}^{0} (2z)^{-n}
= \sum_{m=0}^{\infty} (2z)^{m}
= \frac{1}{1-2z},\qquad |2z|<1 \;\Rightarrow\; |z|<\frac{1}{2}.
\]

\textbf{(4) $\delta(n)$}：单位样值序列的Z变换恒为常数：
\[
\mathcal{Z}[\delta(n)] = 1,\qquad 0 \le |z| \le \infty.
\]

\textbf{(5) $\delta(n-1)$}：利用时移性质 $\delta(n-n_0)\xleftrightarrow{Z} z^{-n_0}$（全平面或除原点/无穷远），得
\[
\mathcal{Z}[\delta(n-1)] = z^{-1},\qquad 0 < |z| \le \infty.
\]

\textbf{(6) $2^{-n}[u(n)-u(n-10)]$}：\textbf{\textcolor{myred}{此为有限长序列（$n=0,1,\ldots,9$），收敛域为整个 $z$ 平面}}。利用有限项等比级数公式：
\[
\mathcal{Z}\big[2^{-n}(u(n)-u(n-10))\big]
= \sum_{n=0}^{9} \big(\tfrac{1}{2}z^{-1}\big)^{n}
= \frac{1 - (\frac{1}{2}z^{-1})^{10}}{1 - \frac{1}{2}z^{-1}}
= \frac{1 - 2^{-10}z^{-10}}{1 - \frac{1}{2}z^{-1}},\qquad 0 < |z| \le \infty.
\]

\prob{15}{}{求以下序列的Z变换及其收敛域，并在$z$平面上画出极零点分布图。\\ (1) $x(n)=R_N(n)$，$N=4$ \\ (2) $x(n)=A r^n \cos(\omega_0 n+\varphi) u(n)$，$r=0.9$，$\omega_0=0.5\pi \, \text{rad}$，$\varphi=0.25\pi \, \text{rad}$ \\ (3) $x(n)=\begin{cases} n & 0 \le n \le N \\ 2N-n & N+1 \le n \le 2N \\ 0 & \text{其他} \end{cases}$，式中$N=4$。}

\textbf{解：}

\imgleft[0.3]{images/image-2026-06-15-20-38-13.png}{
\textbf{(1)}  $R_4(n)=u(n)-u(n-4)\;\xleftrightarrow{Z}\;
\frac{1}{1-z^{-1}}-\frac{z^{-4}}{1-z^{-1}}
=\frac{z^{4}-1}{z^{3}(z-1)}.$

\textbf{零点}：由分子 $z^{4}-1=0$ 得 $z_k=e^{\mathrm{j}\frac{2\pi}{4}k}=e^{\mathrm{j}\frac{\pi}{2}k}\;(k=0,1,2,3)$，即 $z=1,\,j,\,-1,\,-j$，均匀分布在单位圆上；

\textbf{极点}：由分母 $z^{3}(z-1)=0$ 得 $z_{1,2,3}=0$ (三阶极点) 及 $z=1$ (一阶极点)。

\textbf{\textcolor{myred}{$z=1$ 处零极点相消}}，实际极点为 $z=0$（三阶）。零极点图及收敛域如图所示。
}

\imgleft[0.3]{images/image-2026-06-15-20-40-48.png}{
\textbf{(2)} 由欧拉公式将余弦分解为复指数：
\[
x(n)=\frac{A}{2}r^{n}\big(e^{\mathrm{j}\omega_{0}n}e^{\mathrm{j}\varphi}+e^{-\mathrm{j}\omega_{0}n}e^{-\mathrm{j}\varphi}\big)u(n).
\]

将 $x(n)$ 视为两个单边复指数序列的线性组合，利用 Z 变换的线性性质直接写出：
\[
X(z)=\frac{A}{2}\Big(\frac{e^{\mathrm{j}\varphi}}{1-r e^{\mathrm{j}\omega_{0}}z^{-1}}+\frac{e^{-\mathrm{j}\varphi}}{1-r e^{-\mathrm{j}\omega_{0}}z^{-1}}\Big),\quad |z|>r.
\]
\textbf{零极点分析：}
\begin{itemize}[leftmargin=*]
  \item \textbf{极点：} $z_{p1}=r e^{\mathrm{j}\omega_{0}},\;z_{p2}=r e^{-\mathrm{j}\omega_{0}}$，位于半径为 $r$ 的圆上，辐角 $\pm\omega_{0}$；
  \item \textbf{零点：} 由分子 $A\big(\cos\varphi-r\cos(\omega_{0}-\varphi)z^{-1}\big)=0$ 解得收敛域为 $|z|>r$，零极点图如图所示。
\end{itemize}
}

\imgleft[0.3]{images/image-2026-06-15-20-49-04.png}{
\textbf{(3)} 注意到：
\[
x(n+1)=y(n)*y(n).
\]

对上式两边取 Z 变换，利用\textbf{时域卷积定理}与\textbf{时移性质}：
\[
\mathcal{Z}[x(n+1)]=zX(z),\qquad
\mathcal{Z}[y(n)*y(n)]=[Y(z)]^{2},
\]
故得
\[
zX(z)=[Y(z)]^{2}\;\Longrightarrow\;X(z)=z^{-1}[Y(z)]^{2}.
\]

代入即得
\[
X(z)=z^{-1}\Bigl[\frac{z^{4}-1}{z^{3}(z-1)}\Bigr]^{2}
      =\frac{1}{z^{7}}\Bigl[\frac{z^{4}-1}{z-1}\Bigr]^{2}.
\]

\textbf{零极点与收敛域分析：}
\begin{itemize}[leftmargin=*]
  \item \textbf{极点：} 分母 $z^{7}(z-1)^{2}=0$ 给出 $z=0$（七阶极点）与 $z=1$（二阶极点）；
  \item \textbf{零点：} 分子 $(z^{4}-1)^{2}=0$ 给出 $z_{k}=e^{\mathrm{j}\frac{\pi}{2}k}\;(k=0,1,2,3)$，均为二阶零点，其中 $z=1$ 处零点与同位置极点\textbf{精确对消}，故实际极点仅为 $z=0$（七阶）；
  \item \textbf{收敛域：} 序列为有限长（$n=0,1,\dots,2N$），故收敛域为 $0<|z|\leqslant\infty$，包含整个 $z$ 平面（可能除去原点）。
\end{itemize}
}

\prob{16}{}{已知$X(z)=\frac{3}{1-\frac{1}{2}z^{-1}}+\frac{2}{1-2z^{-1}}$，求出对应$X(z)$的各种可能的序列表达式。}

\textbf{解：}

根据两个极点，分为\textbf{全左序列}:
$$x(n)=-\operatorname{Res}[F(z), 0.5]-\operatorname{Res}[F(z), 2] 
$$

$$=-\left.\frac{(5z-7)z^{n}}{(z-0.5)(z-2)}(z-0.5)\right|_{z=0.5}-\left.\frac{(5z-7)z^{n}}{(z-0.5)(z-2)}(z-2)\right|_{z=2}$$

$$=-\left[3\cdot\left(\frac{1}{2}\right)^{n}+2\cdot 2^{n}\right]u(-n-1)$$

\textbf{双边序列：}

$$x(n)=3\cdot\left(\frac{1}{2}\right)^{n}u(n)-2\cdot 2^{n}u(-n-1)$$ 

\textbf{全右序列：}

$$x(n)=\operatorname{Res}[F(z), 0.5]+\operatorname{Res}[F(z), 2]=3\cdot\left(\frac{1}{2}\right)^{n}+2\cdot 2^{n}$$

\prob{17}{}{已知$x(n)=a^n u(n)$，$0<a<1$。分别求：\\ (1) $x(n)$的Z变换；\\ (2) $n x(n)$的Z变换；\\ (3) $a^{-n} u(-n)$的Z变换。}

\textbf{解：}

(1) $X(z)=\mathrm{ZT}[a^{n}u(n)]=\sum\limits_{n=-\infty}^{\infty}a^{n}u(n)z^{-n}=\frac{1}{1-az^{-1}}\quad|z|>a$

(2) $\mathrm{ZT}[nx(n)]=-z\frac{\mathrm{d}}{\mathrm{d}z}X(z)=\frac{-az^{-2}}{(1-az^{-1})^{2}}\quad|z|>a$

(3) $\mathrm{ZT}[a^{-n}u(-n)]=\sum\limits_{n=0}^{-\infty}a^{-n}z^{-n}=\sum\limits_{n=0}^{\infty}a^{n}z^{n}=\frac{1}{1-az}\quad|z|<|a|^{-1}$

\prob{18}{}{已知$X(z)=\frac{-3z^{-1}}{2-5z^{-1}+2z^{-2}}$，分别求：\\ (1) 收敛域$0.5<|z|<2$对应的原序列$x(n)$；\\ (2) 收敛域$|z|>2$对应的原序列$x(n)$。}

\textbf{解：}普通的部分分式展开法，看看答案就成。

\prob{19}{}{用部分分式法求以下$X(z)$的反变换：\\ (1) $X(z)=\frac{1-\frac{1}{3}z^{-1}}{2-5z^{-1}+2z^{-2}}$，$\vert z\vert>\frac{1}{2}$；\\ (2) $X(z)=\frac{1-2z^{-1}}{1-\frac{1}{4}z^{-2}}$，$\vert z\vert<\frac{1}{2}$。}

\textbf{解：}普通的部分分式展开法，看看答案就成。

\prob{21}{}{用Z变换法解下列差分方程：\\ (1) $y(n)-0.9y(n-1)=0.05u(n)$，$y(n)=0$，$n\le-1$；\\ (2) $y(n)-0.9y(n-1)=0.05u(n)$，$y(-1)=1$，$y(n)=0$，$n<-1$；\\ (3) $y(n)-0.8y(n-1)-0.15y(n-2)=\delta(n)$，$y(-1)=0.2$，$y(-2)=0.5$，$y(n)=0$，$n\le-3$。}



\prob{22}{}{设线性时不变系统的系统函数$H(z)$为$H(z)=\frac{1-a^{-1}z^{-1}}{1-az^{-1}}$（$a$为实数）。\\ (1) 在$z$平面上用几何法证明该系统是全通网络，即$\vert H(e^{j\omega})\vert=$常数；\\ (2) 参数$a$如何取值，才能使系统因果稳定？画出其极零点分布及收敛域。}

\prob{23}{}{设系统由下面差分方程描述：$y(n)=y(n-1)+y(n-2)+x(n-1)$。\\ (1) 求系统的系统函数$H(z)$，并画出极零点分布图；\\ (2) 限定系统是因果的，写出$H(z)$的收敛域，并求出其单位脉冲响应$h(n)$；\\ (3) 限定系统是稳定性的，写出$H(z)$的收敛域，并求出其单位脉冲响应$h(n)$。}

\prob{24}{}{已知线性因果网络用下面差分方程描述：$y(n)=0.9y(n-1)+x(n)+0.9x(n-1)$。\\ (1) 求网络的系统函数$H(z)$及单位脉冲响应$h(n)$；\\ (2) 写出网络频率响应函数$H(e^{j\omega})$的表达式，并定性画出其幅频特性曲线；\\ (3) 设输入$x(n)=e^{j\omega_0 n}$，求输出$y(n)$。}

\prob{25}{}{已知网络的输入和单位脉冲响应分别为$x(n)=a^n u(n)$，$h(n)=b^n u(n)$（$0<a<1$，$0<b<1$）。\\ (1) 试用卷积法求网络输出$y(n)$；\\ (2) 试用ZT法求网络输出$y(n)$。}

\prob{27}{}{如果$x_1(n)$和$x_2(n)$是两个不同的因果稳定实序列，求证：$\frac{1}{2\pi}\int_{-\pi}^{\pi}X_1(e^{j\omega})X_2(e^{j\omega})d\omega=\left[\frac{1}{2\pi}\int_{-\pi}^{\pi}X_1(e^{j\omega})d\omega\right]\left[\frac{1}{2\pi}\int_{-\pi}^{\pi}X_2(e^{j\omega})d\omega\right]$，式中$X_1(e^{j\omega})$和$X_2(e^{j\omega})$分别表示$x_1(n)$和$x_2(n)$的傅里叶变换。}

\prob{28}{}{若序列$h(n)$是因果序列，其傅里叶变换的实部如下式：$H_{\text{R}}(e^{j\omega})=\frac{1-a\cos\omega}{1+a^2-2a\cos\omega}$（$\vert a\vert<1$），求序列$h(n)$及其傅里叶变换$H(e^{j\omega})$。}

\prob{29}{}{若序列$h(n)$是因果序列，$h(0)=1$，其傅里叶变换的虚部为$H_{\text{I}}(e^{j\omega})=\frac{-a\sin\omega}{1+a^2-2a\cos\omega}$（$\vert a\vert<1$），求序列$h(n)$及其傅里叶变换$H(e^{j\omega})$。}

\section{第三章\ DFT与 FFT}

\prob{1}{}{计算以下序列的 $N$ 点 DFT，在变换区间 $0 \leqslant n \leqslant N-1$ 内，序列定义为\newline (1) $x(n)=1$\newline (2) $x(n)=\delta(n)$\newline (3) $x(n)=\delta(n-n_0) \quad 0 < n_0 < N$\newline (4) $x(n)=R_m(n) \quad 0 < m < N$\newline (5) $x(n)=\mathrm{e}^{\mathrm{j}\frac{2\pi}{N}mn} \quad 0 < m < N$\newline (6) $x(n)=\cos\left(\frac{2\pi}{N}mn\right) \quad 0 < m < N$\newline (7) $x(n)=\mathrm{e}^{\mathrm{j}\omega_0 n}R_N(n)$\newline (8) $x(n)=\sin(\omega_0 n)\cdot R_N(n)$\newline (9) $x(n)=\cos(\omega_0 n)\cdot R_N(N)$\newline (10) $x(n)=nR_N(n)$}

\prob{2}{}{已知下列 $X(k)$，求 $x(n)=\mathrm{IDFT}[X(k)]$：\newline (1) $X(k)=\begin{cases} \dfrac{N}{2}\mathrm{e}^{\mathrm{j}\theta} & k=m \\[8pt] \dfrac{N}{2}\mathrm{e}^{-\mathrm{j}\theta} & k=N-m \\[8pt] 0 & \text{其他 }k \end{cases}$\newline (2) $X(k)=\begin{cases} -\dfrac{N}{2}\mathrm{e}^{\mathrm{j}\theta} & k=m \\[8pt] \dfrac{N}{2}\mathrm{j}\mathrm{e}^{-\mathrm{j}\theta} & k=N-m \\[8pt] 0 & \text{其他 }k \end{cases}$\newline 其中，$m$ 为正整数，$0 < m < N/2$。}

\prob{3}{}{已知长度为 $N=10$ 的两个有限长序列：\newline $x_1(n)=\begin{cases} 1 & 0 \leqslant n \leqslant 4 \\ 0 & 5 \leqslant n \leqslant 9 \end{cases}$\newline $x_2(n)=\begin{cases} 1 & 0 \leqslant n \leqslant 4 \\ -1 & 5 \leqslant n \leqslant 9 \end{cases}$\newline 作图表示 $x_1(n)$、$x_2(n)$ 和 $y(n)=x_1(n)\circledast x_2(n)$，循环卷积区间长度 $L=10$。}

\prob{4}{}{证明 DFT 的对称定理，即假设 $X(k)=\mathrm{DFT}[x(n)]$，证明\newline $\mathrm{DFT}[X(n)]=Nx(N-k)$}

\prob{5}{}{如果 $X(k)=\mathrm{DFT}[x(n)]$，证明 DFT 的初值定理\newline $x(0)=\dfrac{1}{N}\sum\limits_{k=0}^{N-1}X(k)$}

\prob{6}{}{设 $x(n)$ 的长度为 $N$，且\newline $X(k)=\mathrm{DFT}[x(n)] \quad 0 \leqslant k \leqslant N-1$\newline 令\newline $h(n)=x((n))_N R_{mN}(n) \quad m\text{ 为自然数}$\newline $H(k)=\mathrm{DFT}[h(n)]_{mN} \quad 0 \leqslant k \leqslant mN-1$\newline 求 $H(k)$ 与 $X(k)$ 的关系式。}

\prob{7}{}{证明：若 $x(n)$ 为实序列，$X(k)=\mathrm{DFT}[x(n)]_N$，则 $X(k)$ 为共轭对称序列，即 $X(k)=X^*(N-k)$；若 $x(n)$ 实偶对称，即 $x(n)=x(N-n)$，则 $X(k)$ 也实偶对称；若 $x(n)$ 实奇对称，即 $x(n)=-x(N-n)$，则 $X(k)$ 为纯虚函数并奇对称。}

\prob{8}{}{证明频域循环移位性质：设 $X(k)=\mathrm{DFT}[x(n)]$，$Y(k)=\mathrm{DFT}[y(n)]$，如果 $Y(k)=X((k+l))_N R_N(k)$，则\newline $y(n)=\mathrm{IDFT}[Y(k)]=W_N^{ln}x(n)$}

\prob{9}{}{已知 $x(n)$ 长度为 $N$，$X(k)=\mathrm{DFT}[x(n)]$，\newline $y(n)=\begin{cases} x(n) & 0 \leqslant n \leqslant N-1 \\ 0 & N \leqslant n \leqslant mN-1 \end{cases}$，$m$ 为自然数\newline $Y(k)=\mathrm{DFT}[y(n)]_{mN} \quad 0 \leqslant k \leqslant mN-1$\newline 求 $Y(k)$ 与 $X(k)$ 的关系式。}

\prob{10}{}{证明离散相关定理。若\newline $X(k)=X_1^*(k)X_2(k)$\newline 则\newline $x(n)=\mathrm{IDFT}[X(k)]=\sum\limits_{l=0}^{N-1}x_1^*(l)x_2((l+n))_N R_N(n)$}



\prob{12}{}{已知 $f(n)=x(n)+\mathrm{j}y(n)$，$x(n)$ 与 $y(n)$ 均为长度为 $N$ 的实序列。设\newline $F(k)=\mathrm{DFT}[f(n)]_N \quad 0 \leqslant k \leqslant N-1$\newline (1) \quad $F(k)=\dfrac{1-a^N}{1-aW_N^k}+\mathrm{j}\,\dfrac{1-b^N}{1-bW_N^k} \quad a,b\text{ 为实数}$\newline (2) \quad $F(k)=1+\mathrm{j}N$\newline 试求 $X(k)=\mathrm{DFT}[x(n)]_N$，$Y(k)=\mathrm{DFT}[y(n)]_N$ 以及 $x(n)$ 和 $y(n)$。}

\prob{13}{}{已知序列 $x(n)=a^n u(n)$，$0<a<1$，对 $x(n)$ 的 Z 变换 $X(z)$ 在单位圆上等间隔采样 $N$ 点，采样序列为\newline $X(k)=x(z)\big|_{z=W_N^{-k}} \quad k=0,1,\cdots,N-1$\newline 求有限长序列 $\mathrm{IDFT}[X(k)]_N$。}

\prob{14}{}{两个有限长序列 $x(n)$ 和 $y(n)$ 的零值区间为\newline $x(n)=0 \quad n<0,\,8\leqslant n$\newline $y(n)=0 \quad n<0,\,20\leqslant n$\newline 对每个序列作 20 点 DFT，即\newline $X(k)=\mathrm{DFT}[x(n)] \quad k=0,1,\cdots,19$\newline $Y(k)=\mathrm{DFT}[y(n)] \quad k=0,1,\cdots,19$\newline 如果\newline $F(k)=X(k)\cdot Y(k) \quad k=0,1,\cdots,19$\newline $f(n)=\mathrm{IDFT}[F(k)] \quad k=0,1,\cdots,19$\newline 试问在哪些点上 $f(n)$ 与 $x(n)*y(n)$ 值相等，为什么？}

\prob{15}{}{已知实序列 $x(n)$ 的 8 点 DFT 的前 5 个值为 $0.25$，$0.125-\mathrm{j}0.3018$，$0$，$0.125-\mathrm{j}0.0518$，$0$。\newline (1) 求 $X(k)$ 的其余 3 点的值；\newline (2) $x_1(n)=\sum\limits_{m=-\infty}^{+\infty}x(n+5+8m)R_8(n)$，求 $X_1(k)=\mathrm{DFT}[x_1(n)]_8$；\newline (3) $x_2(n)=x(n)\mathrm{e}^{\mathrm{j}\pi n/4}$，求 $x_2(k)=\mathrm{DFT}[x_2(n)]_8$。}

\prob{16}{}{
已知 $x(n)$、$x_1(n)$、$x_2(n)$ 分别如题 16 图 (a)、(b)、(c) 所示，$X(k)=\mathrm{DFT}[x(n)]_8$。
求 $X_1(k)=\mathrm{DFT}[x_1(n)]_8$ 和 $X_2(k)=\mathrm{DFT}[x_2(n)]_8$（用 $X(k)$ 表示）。

\medskip
\begin{center}
\begin{tikzpicture}[scale=0.55]
  \draw[->] (-0.5,0) -- (8.5,0) node[right] {$n$};
  \node[above] at (0,0.7) {\footnotesize (a) $x(n)$};
  \foreach \x/\y in {0/1,1/1,2/1,3/1,4/1} {\fill (\x,0) circle (1.5pt); \draw (\x,0) -- (\x,\y);}
  \foreach \x in {5,6,7} {\fill (\x,0) circle (1.5pt);}
  \node[below] at (7,-0.2) {\footnotesize $7$};
\end{tikzpicture}
\hfill
\begin{tikzpicture}[scale=0.55]
  \draw[->] (-0.5,0) -- (8.5,0) node[right] {$n$};
  \node[above] at (0,0.7) {\footnotesize (b) $x_1(n)$};
  \foreach \x/\y in {0/1,4/1,5/1,6/1,7/1} {\fill (\x,0) circle (1.5pt); \draw (\x,0) -- (\x,\y);}
  \foreach \x in {1,2,3} {\fill (\x,0) circle (1.5pt);}
  \node[below] at (7,-0.2) {\footnotesize $7$};
\end{tikzpicture}
\hfill
\begin{tikzpicture}[scale=0.55]
  \draw[->] (-0.5,0) -- (8.5,0) node[right] {$n$};
  \node[above] at (0,0.7) {\footnotesize (c) $x_2(n)$};
  \foreach \x/\y in {1/1,2/1,3/1,4/1,5/1} {\fill (\x,0) circle (1.5pt); \draw (\x,0) -- (\x,\y);}
  \foreach \x in {0,6,7} {\fill (\x,0) circle (1.5pt);}
  \node[below] at (7,-0.2) {\footnotesize $7$};
\end{tikzpicture}
\end{center}
\centerline{题 16 图}
}

\prob{17}{}{设 $x(n)$ 是长度为 $N$ 的因果序列，且\newline $X(\mathrm{e}^{\mathrm{j}\omega})=\mathrm{FT}[x(n)]$，$y(n)=\left[\sum\limits_{m=-\infty}^{\infty}x(n+mM)\right]R_M(n)$，$Y(k)=\mathrm{DFT}[y(n)]_M$\newline 试确定 $Y(k)$ 与 $X(\mathrm{e}^{\mathrm{j}\omega})$ 的关系式。}

\prob{18}{}{用微处理机对实数序列作谱分析，要求谱分辨率 $F\leqslant 50$ Hz，信号最高频率为 $1$ kHz，试确定以下各参数：\newline (1) 最小记录时间 $T_{p\min}$；\newline (2) 最大取样间隔 $T_{\max}$；\newline (3) 最少采样点数 $N_{\min}$；\newline (4) 在频带宽度不变的情况下，使频率分辨率提高 1 倍(即 $F$ 缩小一半)的 $N$ 值。}

\prob{19}{}{已知调幅信号的载波频率 $f_c=1$ kHz，调制信号频率 $f_m=100$ Hz，用 FFT 对其进行谱分析，试求：\newline (1) 最小记录时间 $T_{p\min}$；\newline (2) 最低采样频率 $f_{s\min}$；\newline (3) 最少采样点数 $N_{\min}$。}

\comment{3}{}{以下题目来自 GXQ 教材第四章「快速傅里叶变换」}

\prob{1}{}{如果某通用单片计算机的速度为平均每次复数乘需要 $4\,\mu\text{s}$，每次复数加需要 $1\,\mu\text{s}$，用来计算 $N=1024$ 点 DFT，问直接计算需要多少时间。用 FFT 计算呢？照这样计算，用 FFT 进行快速卷积来对信号进行处理时，估计可实现实时处理的信号最高频率。}

\prob{2}{}{如果将通用单片机换成数字信号处理专用单片机 TMS320 系列，计算复数乘和复数加各需要 $10\,\text{ns}$。请重复做上题。}

\prob{3}{}{已知 $X(k)$ 和 $Y(k)$ 是两个 $N$ 点实序列 $x(n)$ 和 $y(n)$ 的 DFT，希望从 $X(k)$ 和 $Y(k)$ 求 $x(n)$ 和 $y(n)$，为提高运算效率，试设计用一次 $N$ 点 IFFT 来完成的算法。}

\prob{4}{}{设 $x(n)$ 是长度为 $2N$ 的有限长实序列，$X(k)$ 为 $x(n)$ 的 $2N$ 点 DFT。\newline (1) 试设计用一次 $N$ 点 FFT 完成计算 $X(k)$ 的高效算法；\newline (2) 若已知 $X(k)$，试设计用一次 $N$ 点 IFFT 实现求 $X(k)$ 的 $2N$ 点 IDFT 运算。}

\prob{5}{}{分别画出 $16$ 点基 2 DIT-FFT 和 DIF-FFT 运算流图，并计算其复数乘次数，如果考虑三类碟形的乘法计算，试计算复乘次数。}

\section{第四章\ 时域离散系统的算法结构}

\prob{1}{}{已知系统用下面差分方程描述：\newline $y(n)=\frac{3}{4}y(n-1)-\frac{1}{8}y(n-2)+x(n)+\frac{1}{3}x(n-1)$\newline 试分别画出系统的直接型、级联型和并联型结构。式中 $x(n)$ 和 $y(n)$ 分别表示系统的输入和输出信号。}

\textbf{解：}

（1）

\imgleft[0.3]{images/image-2026-06-18-15-20-14.png}{

原式移项得：
$y(n) - \frac{3}{4}y(n-1) + \frac{1}{8}y(n-2) = x(n) + \frac{1}{3}x(n-1)$

Z变换，得到:
$Y(z) - \frac{3}{4}Y(z)z^{-1} + \frac{1}{8}Y(z)z^{-2} = X(z) + \frac{1}{3}X(z)z^{-1}$
$$H(z) = \frac{1 + \frac{1}{3}z^{-1}}{1 - \frac{3}{4}z^{-1} + \frac{1}{8}z^{-2}}$$

}

\textbf{（2）级联型结构：}

\imgleft[0.3]{images/image-2026-06-18-15-32-00.png}{
将$H(z)$的分母进行因式分解：

$H(z) = \frac{1 + \frac{1}{3}z^{-1}}{1 - \frac{3}{4}z^{-1} + \frac{1}{8}z^{-2}} = \frac{1 + \frac{1}{3}z^{-1}}{\left(1 - \frac{1}{2}z^{-1}\right)\left(1 - \frac{1}{4}z^{-1}\right)}$

按照上式可以有两种级联型结构：

$H(z) = \frac{1 + \frac{1}{3}z^{-1}}{1 - \frac{1}{2}z^{-1}} \cdot \frac{1}{1 - \frac{1}{4}z^{-1}}$

$H(z) = \frac{1}{1 - \frac{1}{2}z^{-1}} \cdot \frac{1 + \frac{1}{3}z^{-1}}{1 - \frac{1}{4}z^{-1}}$
}

\textbf{（3）并联型结构：}

\imgleft[0.3]{images/image-2026-06-18-15-45-19.png}{
$\frac{H(z)}{z} = \frac{\frac{10}{3}}{z - \frac{1}{2}} - \frac{\frac{7}{3}}{z - \frac{1}{4}}$

$H(z) = \frac{\frac{10}{3}z}{z - \frac{1}{2}} - \frac{\frac{7}{3}z}{z - \frac{1}{4}} = \frac{\frac{10}{3}}{1 - \frac{1}{2}z^{-1}} + \frac{-\frac{7}{3}}{1 - \frac{1}{4}z^{-1}}$
}

\prob{2}{}{设数字滤波器的差分方程为\newline $y(n)=x(n)+x(n-1)+\frac{1}{3}y(n-1)+\frac{1}{4}y(n-2)$\newline 试画出系统的直接型结构。}

\imgleft[0.3]{images/image-2026-06-18-15-46-50.png}{
\textbf{解：}

由差分方程得到滤波器的系统函数为  

$H(z)=\frac{1+z^{-1}}{1-\frac{1}{3}z^{-1}-\frac{1}{4}z^{-2}}$  

画出其直接型结构如题 2 解图所示。
}

\prob{3}{}{设系统的差分方程为\newline $y(n)=(a+b)y(n-1)-aby(n-2)+x(n-2)+(a+b)x(n-1)+abx(n)$\newline 式中，$|a|<1$，$|b|<1$，$x(n)$ 和 $y(n)$ 分别表示系统的输入和输出信号，试画出系统的直接型和级联型结构。}

\textbf{解：}

\textbf{（1）直接型结构：}

\imgleft[0.3]{images/image-2026-06-18-15-53-20.png}{
将差分方程进行Z变换，得到  

$Y(z) = (a + b)Y(z)z^{-1} - abY(z)z^{-2} + X(z)z^{-2} - (a + b)X(z)z^{-1} + ab$  

$H(z) = \frac{Y(z)}{X(z)} = \frac{ab - (a + b)z^{-1} + z^{-2}}{1 - (a + b)z^{-1} - abz^{-2}}$
}

\textbf{（2）级联型结构：}

\imgleft[0.3]{images/image-2026-06-18-19-59-06.png}{
$H(z) = \frac{(a - z^{-1})(b - z^{-1})}{(1 - az^{-1})(1 - bz^{-1})} = H_1(z)H_2(z)$  

\textbf{a:}$H_1(z) = \frac{z^{-1} - a}{1 - az^{-1}}, H_2(z) = \frac{z^{-1} - b}{1 - bz^{-1}}$  
  
\textbf{b:}$H_1(z) = \frac{z^{-1} - a}{1 - bz^{-1}}, H_2(z) = \frac{z^{-1} - b}{1 - az^{-1}}$
}

\prob{4}{}{设系统的系统函数为\newline $H(z)=4\frac{(1+z^{-1})(1-1.414z^{-1}+z^{-2})}{(1-0.5z^{-1})(1+0.9z^{-1}+0.81z^{-2})}$\newline 试画出各种可能的级联型结构，并指出哪一种最好。}

\textbf{解：}

\imgleft[0.3]{images/image-2026-06-18-20-03-51.png}{
根据系统函数表达式易得有两种级联方式：
 
$H_1(z) = \frac{4(1 + z^{-1})}{1 - 0.5z^{-1}}, H_2(z) = \frac{1 - 1.414z^{-1} + z^{-2}}{1 + 0.9z^{-1} + 0.81z^{-2}}$  

$H_1(z) = \frac{1 - 1.414z^{-1} + z^{-2}}{1 - 0.5z^{-1}}, H_2(z) = \frac{4(1 + z^{-1})}{1 + 0.9z^{-1} + 0.81z^{-2}}$

\textbf{\textcolor{myred}{第一种更好：将需要延时两级的做成一个子系统，使用的延时器总数最少，也更方便。}}
}

\prob{5}{}{
  题 5 图中画出了四个系统，试用各子系统的单位脉冲响应分别表示各总系统的
  单位脉冲响应，并求其总系统函数。
  \img{images/image-2026-06-17-11-37-38.png}
}

\textbf{解：}

（1）$h(n)=h_{1}(n)*h_{2}(n)*h_{3}(n)$, $H(z)=H_{1}(z)H_{2}(z)H_{3}(z)$

（2）$h(n)=h_{1}(n)+h_{2}(n)+h_{3}(n)$, $H(z)=H_{1}(z)+H_{2}(z)+H_{3}(z)$

（3）$h(n)=h_{1}(n)*h_{2}(n)+h_{3}(n)$, $H(z)=H_{1}(z)\cdot H_{2}(z)+H_{3}(z)$

（4）$h(n)=h_{1}(n)*[h_{2}(n)+h_{3}(n)*h_{4}(n)]+h_{5}(n)$

$\quad =h_{1}(n)*h_{2}(n)+h_{1}(n)*h_{3}(n)*h_{4}(n)+h_{5}(n)$

$H(z)=H_{1}(z)H_{2}(z)+H_{1}(z)H_{3}(z)H_{4}(z)+H_{5}(z)$

\prob{6}{}{
  题 6 图中画出了10 种不同的流图，分别写出它们的系统函数以及差分方程。
  \img{images/image-2026-06-17-11-39-27.png}
}

\textbf{解：}

由于题干里重复的类型太多了，这里只详细探讨（c）（h）问，其余小问考前看一眼即可。

\textbf{（c）：}

典型的 FIR 系统结构流图，但是就是一个简单类似于并联的结构，可得系统函数：$H(z)=a+bz^{-1}+cz^{-2}$

\textbf{（h）：}

\textbf{\textcolor{mygreen}{还是没怎么看懂（2026-06-19D）}}。

\prob{7}{}{假设滤波器的单位脉冲响应为\newline $h(n)=a^n u(n) \quad 0<a<1$\newline 求出滤波器的系统函数，并画出它的直接型结构。}

\textbf{解：}

\imgleft[0.3]{images/image-2026-06-19-03-48-22.png}{
由 Z 变换，得到该滤波器的系统函数为：$H(z)=\frac{1}{1-az^{-1}}$
}

\prob{8}{}{已知系统的单位脉冲响应为\newline $h(n)=\delta(n)+2\delta(n-1)+0.3\delta(n-2)+2.5\delta(n-3)+0.5\delta(n-5)$\newline 试写出系统的系统函数，并画出它的直接型结构。}

\textbf{解：}

典型的 FIR 系统的冲激响应结构，由 Z 变换，系统函数为：$H(z) = 1 + 2z^{-1} + 0.3z^{-2} + 2.5z^{-3} + 0.5z^{-5}$，由此可绘制直接型信号流图。

\img{images/image-2026-06-19-03-53-49.png}


\prob{9}{\textbf{\textcolor{myred}{线性相位结构母题}}}{已知 FIR 滤波器的系统函数为\newline $H(z)=\frac{1}{10}(1+0.9z^{-1}+2.1z^{-2}+0.9z^{-3}+z^{-4})$\newline 试画出该滤波器的直接型结构和线性相位结构。}

\textbf{解：}

主要关注本题的后半部分（奇数项冲击序列的线性相位结构）

\imgleft[0.3]{images/image-2026-06-19-04-02-39.png}{
绘制这个图的理解比较抽象（类似于电路原理的那种等效电路概念），这里暂时无法使用语言来描述这个过程，留一个链接（正好是本题的讲解视频，仅供参考）：\textbf{\textcolor{myblue}{\href{https://www.bilibili.com/video/BV1do5Fz5EMf?p=71}{讲解视频}}}

但总的思路就是理解信号“流“的概念，然后就能大致明白为什么会使用线性相位结构这种“复用“结构可以简化实现了。
}

\prob{10}{}{已知 FIR 滤波器的单位脉冲响应为：\newline (1) $N=6$\newline $h(0)=h(5)=15$\newline $h(1)=h(4)=2$\newline $h(2)=h(3)=3$\newline (2) $N=7$\newline $h(0)=-h(6)=3$\newline $h(1)=-h(5)=-2$\newline $h(2)=-h(4)=1$\newline $h(3)=0$\newline 试画出它们的线性相位型结构图，并分别说明它们的\textbf{\textcolor{myred}{幅度特性、相位特性}}各有什么特点。}

\textbf{解：}

有 T9 的经验之后，线性相位结构图就好画了，这里不再展示指导书上的图了（因为指导书疑似出现了一些箭头标反的操作，这里不再展示谨防误导）


\textbf{幅度、相位特性：}见 cheatsheet 上整理的第一类、第二类 FIR 系统的特性总结，直接抄就行。

\prob{11}{}{已知 FIR 滤波器的 16 个频率采样值为：\newline $H(0)=12$，\qquad $H(3)\sim H(13)=0$\newline $H(1)=-3-\mathrm{j}\sqrt{3}$，\qquad $H(14)=1-\mathrm{j}$\newline $H(2)=1+\mathrm{j}$，\qquad $H(15)=-3+\mathrm{j}\sqrt{3}$\newline 试画出其频率采样结构，选择 $r=1$，可以用复数乘法器。}

\textbf{解：}

\imgleft[0.4]{images/image-2026-06-19-05-04-56.png}{
核心就是\textbf{\textcolor{myred}{把题干信息代入到频率采样格式的系统函数公式中}}，对于本题，即：

$$H(z) = \frac{1 - z^{-N}}{N} \sum_{k=0}^{N-1} \frac{H(k)}{1 - W_N^{-k} z^{-1}} \quad N = 16$$

拆分为“级联+并联“的形式：

\entry{级联流图系数}{即$\frac{1}{N}$}

\entry{并联流图系数}{即$W_N^{-k}$，以及对应的频率采样值$H(k)$}
}

\prob{12}{\textbf{\textcolor{myred}{修正频率采样结构（ban 复数乘法器）}}}{已知 FIR 滤波器系统函数在单位圆上 16 个等间隔采样点为：\newline $H(0)=12$\qquad $H(3)\sim H(13)=0$\newline $H(1)=-3-\mathrm{j}\sqrt{3}$\qquad $H(14)=1-\mathrm{j}$\newline $H(2)=1+\mathrm{j}$\qquad $H(15)=-3+\mathrm{j}\sqrt{3}$\newline 试画出它的频率采样结构，取修正半径 $r=0.9$，要求用实数乘法器。}

\textbf{解：}

主要考察修正频率采样机构表达式，剩下的就交给卡西欧了：

$$H(z) = \frac{1 - r^{16} z^{-16}}{16} \sum_{k=0}^{16} \frac{H(k)}{1 - r W_{16}^{-k} z^{-1}}$$

$$= \frac{1}{16} (1 - 0.1853 z^{-16}) \left[ \frac{H(0)}{1 - 0.9 z^{-1}} + \left( \frac{H(1)}{1 - 0.9 W_{16}^{-1} z^{-1}} + \frac{H(15)}{1 - 0.9 W_{16}^{-15} z^{-1}} \right) \right.$$

$$\left. + \left( \frac{H(2)}{1 - 0.9 W_{16}^{-2} z^{-1}} + \frac{H(14)}{1 - 0.9 W_{16}^{-14} z^{-1}} \right) \right]$$

$$= \frac{1}{16} (1 - 0.1853 z^{-16}) \left[ \frac{12}{1 - 0.9 z^{-1}} + \left( \frac{-6 - 6.182 z^{-1}}{1 - 1.663 z^{-1} + 0.81 z^{-2}} \right. \right.$$

$$\left. \left. + \frac{2 - 2.5456 z^{-1}}{1 - 1.2728 z^{-1} + 0.81 z^{-2}} \right) \right]$$

\prob{13}{}{已知 FIR 滤波器的单位脉冲响应为\newline $h(n)=\delta(n)-\delta(n-1)+\delta(n-4)$\newline 试用频率采样结构实现该滤波器。设采样点数 $N=5$，要求画出频率采样网络结构，写出滤波器参数的计算公式。}

\textbf{解：}

我们已知频率采样形式下的系统函数表达式：

$$H(z) = \frac{1 - z^{-N}}{N} \sum_{k=0}^{N-1} \frac{H(k)}{1 - W_N^{-k} z^{-1}}$$

现在已知$N = 5$（级联型流图可画），接下来重点就是求解$H(k)$，即对$h(n)$进行$N = 5$点的 DFT，不过答案似乎也懒得算了，直接在系数位置标注$H(0),H(1),H(2),H(3),H(4)$了。

\section{第五章\ IIR数字滤波器的设计}

\prob{1}{}{设计一个巴特沃斯低通滤波器，要求通带截止频率 $f_p=6\,\text{kHz}$，通带最大衰减 $\alpha_p=3\,\text{dB}$，阻带截止频率 $f_s=12\,\text{kHz}$，阻带最小衰减 $\alpha_s=25\,\text{dB}$。求出滤波器归一化系统函数 $G(p)$ 以及实际的 $H_a(s)$。}

\prob{2}{}{设计一个切比雪夫低通滤波器，要求通带截止频率 $f_p=3\,\text{kHz}$，通带最大衰减 $\alpha_p=0.2\,\text{dB}$，阻带截止频率 $f_s=12\,\text{kHz}$，阻带最小衰减 $\alpha_s=50\,\text{dB}$。求出滤波器归一化系统函数 $G(p)$ 和实际的 $H_a(s)$。}

\prob{3}{}{设计一个巴特沃斯高通滤波器，要求其通带截止频率 $f_p=20\,\text{kHz}$，阻带截止频率 $f_s=10\,\text{kHz}$，$f_p$ 处最大衰减为 $3\,\text{dB}$，阻带最小衰减 $\alpha_s=15\,\text{dB}$。求出该高通滤波器的系统函数 $H_a(s)$。}

\prob{4}{}{已知模拟滤波器的系统函数 $H_a(s)$ 如下：\newline (1) $H_a(s)=\dfrac{s+a}{(s+a)^2+b^2}$\newline (2) $H_a(s)=\dfrac{b}{(s+a)^2+b^2}$\newline 式中 $a$、$b$ 为常数，设 $H_a(s)$ 因果稳定，试采用脉冲响应不变法将其转换成数字滤波器 $H(z)$。}

\prob{5}{}{已知模拟滤波器的系统函数如下：\newline (1) $H_a(s)=\dfrac{1}{s^2+s+1}$\newline (2) $H_a(s)=\dfrac{b}{2s^2+3s+1}$\newline 试采用脉冲响应不变法和双线性变换法将其转换为数字滤波器。设 $T=2\,\text{s}$。}

\prob{6}{}{设 $h_a(t)$ 表示一模拟滤波器的单位冲激响应，即\newline $h_a(t) = \begin{cases} e^{-0.9t} & t \geq 0 \\ 0 & t < 0 \end{cases}$\newline 用脉冲响应不变法，将此模拟滤波器转换成数字滤波器（用 $h(n)$ 表示单位脉冲响应，即 $h(n)=h_a(nT)$）。确定系统函数 $H(z)$，并把 $T$ 作为参数，证明：$T$ 为任何值时，数字滤波器是稳定的，并说明数字滤波器近似为低通滤波器还是高通滤波器。}

\prob{7}{}{假设某模拟滤波器 $H_a(s)$ 是一个低通滤波器，又知 $H(z)=H_a(s)\big|_{s=\frac{z-1}{z+1}}$，数字滤波器 $H(z)$ 的通带中心位于下面哪种情况？并说明原因。\newline (1) $\omega=0$（低通）。\newline (2) $\omega=\pi$（高通）。\newline (3) 除 $0$ 或 $\pi$ 以外的某一频率（带通）。}

\prob{9}{}{设计低通数字滤波器，要求通带内频率低于 $0.2\pi$ rad 时，容许幅度误差在 $1$ dB 之内；频率在 $0.3\pi$ 到 $\pi$ 之间的阻带衰减大于 $10$ dB。试采用巴特沃斯型模拟滤波器进行设计，用脉冲响应不变法进行转换，采样间隔 $T=1$ ms。}

\prob{10}{}{要求同题 9，试采用双线性变换法设计数字低通滤波器。}

\prob{11}{}{设计一个数字高通滤波器，要求通带截止频率 $\omega_p=0.8\pi$ rad，通带衰减不大于 3 dB，阻带截止频率 $\omega_s=0.5\pi$ rad，阻带衰减不小于 18 dB，希望采用巴特沃斯型滤波器。}

\prob{12}{}{设计一个数字带通滤波器，通带范围为 $0.25\pi$ rad 到 $0.45\pi$ rad，通带内最大衰减为 3 dB，$0.15\pi$ rad 以下和 $0.55\pi$ rad 以上为阻带，阻带内最小衰减为 15 dB。试采用巴特沃斯型模拟低通滤波器。}

\section{第六章\ \ FIR数字滤波器的设计}

\prob{1}{}{已知 FIR 滤波器的单位脉冲响应为：
\newline
(1) $h(n)$ 长度 $N=6$
\newline
$h(0)=h(5)=1.5$，$h(1)=h(4)=2$，$h(2)=h(3)=3$
\newline
(2) $h(n)$ 长度 $N=7$
\newline
$h(0)=-h(6)=3$，$h(1)=-h(5)=-2$，$h(2)=-h(4)=1$，$h(3)=0$
\newline
试分别说明它们的幅度特性和相位特性各有什么特点。}

\prob{2}{}{已知第一类线性相位 FIR 滤波器的单位脉冲响应长度为 16，其 16 个频域幅度采样值中的前 9 个为：
\newline
$H_g(0)=12$，$H_g(1)=8.34$，$H_g(2)=3.79$，$H_g(3)\sim H_g(8)=0$
\newline
根据第一类线性相位 FIR 滤波器幅度特性 $H_g(\omega)$ 的特点，求其余 7 个频域幅度采样值。}

\prob{3}{}{设 FIR 滤波器的系统函数为
\newline
$H(z)=\dfrac{1}{10}(1+0.9z^{-1}+2.1z^{-2}+0.9z^{-3}+z^{-4})$
\newline
求出该滤波器的单位脉冲响应 $h(n)$，判断是否具有线性相位，求出其幅度特性函数和相位特性函数。}

\prob{4}{}{用矩形窗设计线性相位低通 FIR 滤波器，要求过渡带宽度不超过 $\pi/8$ rad。希望逼近的理想低通滤波器频率响应函数 $H_d(e^{j\omega})$ 为
\newline
$H_d(e^{j\omega}) = \begin{cases} e^{-j\omega\alpha} & 0 \leq |\omega| \leq \omega_c \\ 0 & \omega_c < |\omega| \leq \pi \end{cases}$
\newline
(1) 求出理想低通滤波器的单位脉冲响应 $h_d(n)$；
\newline
(2) 求出加矩形窗设计的低通 FIR 滤波器的单位脉冲响应 $h(n)$ 表达式，确定 $\alpha$ 与 $N$ 之间的关系；
\newline
(3) 简述 $N$ 取奇数或偶数对滤波特性的影响。}

\prob{5}{}{用矩形窗设计线性相位高通滤波器，要求过渡带宽度不超过 $\pi/10$ rad。希望逼近的理想高通滤波器频率响应函数 $H_d(e^{j\omega})$ 为
\[
H_d(e^{j\omega}) = 
\begin{cases} 
e^{-j\omega\alpha} & \omega_c \leq \omega \leq \pi \\
0 & \text{其他}
\end{cases}
\]
(1) 求出该理想高通的单位脉冲响应 $h_d(n)$；
(2) 求出加矩形窗设计的高通 FIR 滤波器的单位脉冲响应 $h(n)$ 表达式，确定 $\alpha$ 与 $N$ 的关系；
(3) $N$ 的取值有什么限制？为什么？}

\prob{6}{}{理想带通特性为
\[
H_d(e^{j\omega}) = 
\begin{cases} 
e^{-j\omega\alpha} & \omega_c \leq |\omega| \leq \omega_c + B \\
0 & |\omega| < \omega_c,\ \omega_c + B < |\omega| \leq \pi
\end{cases}
\]
(1) 求出该理想带通的单位脉冲响应 $h_d(n)$；
(2) 写出用升余弦窗设计的滤波器的 $h(n)$ 表达式，确定 $\alpha$ 与 $N$ 之间的关系；
(3) 要求过渡带宽度不超过 $\pi/16$ rad。$N$ 的取值是否有限制？为什么？}

\prob{7}{}{试完成下面两题：
(1) 设低通滤波器的单位脉冲响应与频率响应函数分别为 $h(n)$ 和 $H(e^{j\omega})$，另一个滤波器的单位脉冲响应为 $h_1(n)$，它与 $h(n)$ 的关系是 $h_1(n) = (-1)^n h(n)$。试证明滤波器 $h_1(n)$ 是一个高通滤波器。
(2) 设低通滤波器的单位脉冲响应与频率响应函数分别为 $h(n)$ 和 $H(e^{j\omega})$，截止频率为 $\omega_c$，另一个滤波器的单位脉冲响应为 $h_2(n)$，它与 $h(n)$ 的关系是 $h_2(n) = 2h(n)\cos\omega_0 n$，且 $\omega_c < \omega_0 < (\pi - \omega_c)$。试证明滤波器 $h_2(n)$ 是一个带通滤波器。}

\prob{8}{}{题 8 图中 $h_1(n)$ 是偶对称序列，$N=8$，设\newline $H_1(k)=\text{DFT}[h_1(n)] \quad k=0,1,\cdots,N-1$\newline $H_2(k)=\text{DFT}[h_2(n)] \quad k=0,1,\cdots,N-1$\newline (1) 试确定 $H_1(k)$ 与 $H_2(k)$ 的关系式。$|H_1(k)|=|H_2(k)|$ 是否成立？为什么？\newline (2) 用 $h_1(n)$ 和 $h_2(n)$ 分别构成的低通滤波器是否具有线性相位？群延时为多少？}

\prob{9}{}{对下面的每一种滤波器指标，选择满足 FIRDF 设计要求的窗函数类型和长度。\newline (1) 阻带衰减为 $20$ dB，过渡带宽度为 $1$ kHz，采样频率为 $12$ kHz；\newline (2) 阻带衰减为 $50$ dB，过渡带宽度为 $2$ kHz，采样频率为 $20$ kHz；\newline (3) 阻带衰减为 $50$ dB，过渡带宽度为 $500$ Hz，采样频率为 $5$ kHz。}

\prob{10}{}{利用矩形窗、升余弦窗、改进升余弦窗和布莱克曼窗设计线性相位 FIR 低通滤波器。要求希望逼近的理想低通滤波器通带截止频率 $\omega_c=\pi/4$ rad，$N=21$。求出分别对应的单位脉冲响应。\newline (1) 求出分别对应的单位脉冲响应 $h(n)$ 的表达式。}

\prob{11}{}{将技术要求改为设计线性相位高通滤波器，重复题 10。}

\prob{12}{}{利用窗函数(哈明窗)法设计一数字微分器，逼近题 12 图所示的理想微分器特性，并绘出其幅频特性。}

\prob{13}{}{用窗函数法设计一个线性相位低通 FIRDF，要求通带截止频率为 $\pi/4$ rad，过渡带宽度为 $8\pi/51$ rad，阻带最小衰减为 $45$ dB。\newline (1) 选择合适的窗函数及其长度，求出 $h(n)$ 的表达式。\newline (2*) 用 MATLAB 画出损耗函数曲线和相频特性曲线。}

\section{知识点整理寄存器}

\subsection{前言}

傅里叶变换与傅里叶级数的相关概念复习：

\concepttable{
    \conrow{FT}{$X(j\Omega) = \int_{-\infty}^{\infty} x(t)e^{-j\Omega t}dt$，时域连续非周期 $\mapsto$ 频域连续复函数}
    \conrow{IFT}{$x(t) = \frac{1}{2\pi} \int_{-\infty}^{\infty} X(j\Omega)e^{j\Omega t}d\Omega$，归一化系数 $1/2\pi$}
    \conrow{频率变量}{\textbf{\textcolor{myred}{$\Omega$ 为模拟角频率（rad/s）}}，区别于数字角频率 $\omega$}
}

\concepttable{
    \conrow{FS 展开}{$x(t) = \sum_{m=-\infty}^{\infty} C_m e^{jm\omega_0 t}$，基频 $\omega_0 = \frac{2\pi}{T_0}$}
    \conrow{FS 系数}{$C_m = \frac{1}{T_0} \int_{T_0} x(t)e^{-jm\omega_0 t}dt$}
    \conrow{FS vs.\ FT}{\textbf{\textcolor{myred}{时域周期 $\mapsto$ 频域离散（谱线）；非周期 $\mapsto$ 连续谱}}}
}

\concepttable{
    \conrow{时域卷积}{$x(t) * y(t) \overset{\text{FT}}{\leftrightarrow} X(j\Omega) \cdot Y(j\Omega)$}
    \conrow{系统分析}{$y(t)=x(t)*h(t) \;\Rightarrow\; Y(j\Omega)=X(j\Omega)H(j\Omega)$}
    \conrow{频域卷积}{$x(t) \cdot y(t) \overset{\text{FT}}{\leftrightarrow} \frac{1}{2\pi} X(j\Omega) * Y(j\Omega)$}
    \conrow{采样依据}{理想采样 $x(t) \cdot \delta_T(t)$，频域卷积 $\Rightarrow$ 频谱周期延拓}
}

\textbf{\hlblue{傅里叶变换的充要条件}}

\concepttable{
    \conrow{1. 充分条件（狄利克雷）}{$\int_{-\infty}^{\infty}|x(t)|dt<\infty$，绝对可积}
    \conrow{2. 次强条件（能量有限）}{$\int_{-\infty}^{\infty}|x(t)|^2dt<\infty$}
    \conrow{3. 弱条件（功率有限）}{$\lim_{T\to\infty}\frac{1}{2T}\int_{-T}^{T}|x(t)|^2dt<\infty$}
    \conrow{4. 周期信号}{周期信号为功率有限信号的子类；可用FS求$C_m$，也可用广义FT}
}

\compTable{4}{
    \compheader{$x(t)$ & $X(j\Omega)$ & $x(t)$ & $X(j\Omega)$}
    \comprow{$\delta(t)$ & $1$ & $u(t)$ & $\pi\delta(\Omega)+\frac{1}{j\Omega}$}
    \comprow{$1$ & $2\pi\delta(\Omega)$ & $\operatorname{sgn}(t)$ & $\frac{2}{j\Omega}$}
    \comprow{$e^{j\Omega_0 t}$ & $2\pi\delta(\Omega-\Omega_0)$ & $e^{-at}u(t),\,a>0$ & $\frac{1}{a+j\Omega}$}
    \comprow{$\cos(\Omega_0 t)$ & $\pi[\delta(\Omega-\Omega_0)+\delta(\Omega+\Omega_0)]$ & $\operatorname{rect}\!\left(\frac{t}{\tau}\right)$ & $\tau\operatorname{sinc}\!\left(\frac{\Omega\tau}{2}\right)$}
    \comprow{$\sin(\Omega_0 t)$ & $\frac{\pi}{j}[\delta(\Omega-\Omega_0)-\delta(\Omega+\Omega_0)]$ & $\frac{W}{\pi}\operatorname{sinc}(Wt)$ & $\operatorname{rect}\!\left(\frac{\Omega}{2W}\right)$}
    \comprow{$\sum_{n=-\infty}^{\infty}\delta(t-nT)$ & \multicolumn{3}{L}{$\frac{2\pi}{T}\sum_{k=-\infty}^{\infty}\delta\!\left(\Omega-\frac{2\pi k}{T}\right)$}}
}

\textbf{\hlblue{广义傅里叶变换}}

广义FT将$\delta$函数引入频域，使周期信号（不满足绝对可积条件）也能用FT统一分析。

\concepttable{
    \conrow{核心思想}{\textbf{\textcolor{myred}{引入冲激函数 $\delta(\Omega)$}}，将周期信号的频谱表示为频域冲激串，绕过积分发散问题}
    \conrow{数学基础}{$\mathcal{F}\{e^{j\Omega_0 t}\} = 2\pi\delta(\Omega - \Omega_0)$，即复指数信号的频谱为频域单一冲激}
    \conrow{FT对}{$x(t)=e^{j\Omega_0 t} \;\overset{\text{FT}}{\longleftrightarrow}\; X(j\Omega)=2\pi\delta(\Omega-\Omega_0)$}
}

\concepttable{
    \conrow{周期信号广义FT}{由FS展开 $x(t)=\sum_{m=-\infty}^{\infty}C_m e^{jm\omega_0 t}$，两边取广义FT得：}
    \conrow{频谱表达式}{$X(j\Omega)=\sum_{m=-\infty}^{\infty}2\pi C_m\,\delta(\Omega-m\omega_0)$}
    \conrow{物理含义}{时域周期 $\mapsto$ 频域离散冲激串；谱线位于谐波频率 $m\omega_0$，强度为 $2\pi C_m$}
}

\textbf{数学证明}

设 $x_a(t)=A\cos(2\pi F_0 t+\theta)$，另取 $F_k = F_0 + kF_s\ (k \in \mathbb{N})$，则 $x_a'(t)=A\cos(2\pi F_k t+\theta)$。

以 $F_s$ 采样，$t=n/F_s$：
$$x'(n)=A\cos\!\left(2\pi\frac{F_0+kF_s}{F_s}n+\theta\right)=A\cos\!\left(\frac{2\pi nF_0}{F_s}+\theta+2\pi kn\right)$$

由 $\cos(\phi+2\pi kn)=\cos(\phi)$（$k,n\in\mathbb{Z}$），得
$$x'(n)=A\cos\!\left(\frac{2\pi nF_0}{F_s}+\theta\right)=x(n)$$

\concepttable{%
\conrow{有理分式形式}{$X(z) = \frac{P(z)}{Q(z)}$}%
\conrow{零点}{$P(z_0)=0$ 且 $Q(z_0)\neq 0$，记 $z=z_0$，复平面用 $\circ$ 表示。}%
\conrow{极点}{$Q(z_0)=0$，$|X(z_0)|=\infty$，记 $z=z_0$，复平面用 $\times$ 表示。}%
\conrow{结论}{极点处级数必发散，\textbf{\textcolor{myred}{收敛域内部绝不能包含任何极点}}；边界 $R_{x-}$、$R_{x+}$ 由极点半径决定。}%
}

\begin{multicols}{2}
\begin{minipage}{\columnwidth}
\textbf{\hlgreen{椭圆滤波器}}

$|H_a(j\Omega)|^2 = \frac{1}{1 + \varepsilon^2 U_N^2(\Omega/\Omega_p)}$

\textbf{波纹分布：}通带和阻带内都是等波纹。

在相同阶数 $N$ 的情况下，\textbf{陡峭度：椭圆 > 切比雪夫 > 巴特沃斯}。

\end{minipage}
\begin{minipage}{\columnwidth}
\textbf{\hlgreen{贝塞尔滤波器}}

$H_a(s) = \frac{d_0}{B_N(s)}$，分母为贝塞尔多项式

\textbf{幅频特性：}单调变化，无波纹。在相同阶数 $N$ 下，其过渡带比其他滤波器都平缓

但\textbf{相频特性最好，在通带内接近线性相位}。 

\textbf{极点位置：}$\boxed{s_k = \Omega_c e^{j\pi (\frac{1}{2} + \frac{2k+1}{2N})}}$

\end{minipage}
\end{multicols}

% --- 手动换栏命令（如果需要强制换列）---
% \columnbreak 

\end{multicols*}

\end{document}